\ssscp{Distribution of words for ‘tongue’}{15-04-04-02}
\frf{15-09} shows the distribution of known reflexes of
**maya ‘tongue’, and several other etyma
also meaning ‘tongue’ that persist in the region.

The variety of words for tongue is in many ways a reflection
of the complexity and diversity that we find in linguistic Wallacea for other things as well.
Reflexes of **maya are not found in most \tsc{\ili{Bima}-Lembata} languages,
but (unlike **malip) it is found in \ili{Taliabu} (a \tsc{Celebic} language).
Reflexes of **maya are not found in most \tsc{Seram-Tanimbar-Bomberai}
languages, where reflexes of **kəRi are prevalent.
Note in Tanimbar that \ili{Yamdena}
and \ili{Fordata} follow the STB pattern,
while \ili{Selaru}, \ili{Makatian}, and \ili{Seluwasan} follow the pattern found
in other parts of \tsc{Timor-Babar},
consistent with their different subgrouping affiliations.
Reflexes of **maya are not found in \tsc{Aru} languages,
and we note that forms similar to those found in the \tsc{Aru}
subgroup are also found elsewhere.
The sporadic distribution and discontinuities of several of these
forms continue out to Oceanic languages all the way to Polynesia
with Proto-Polynesian *qalelo ‘tongue’.

\begin{figure}[h]\centering
	\includegraphics[width=\columnwidth]{tongue.pdf}
	\fcp{Distribution of words meaning ‘tongue’}{15-09}
\end{figure}

\largerpage
In spite of the claim to be a “replacement innovation” in \citet[740]{Blust2013},
PCEMP **maya ‘tongue’ is not found through the whole CEMP region,
and is not a replacement innovation.
\citet{BlustTrussel2020} curiously assign some of the same forms in the
same “CMP” languages to both PCEMP **maya ‘tongue’
and to \tsc{PAn} *Səma ‘tongue’.
Forms such as \ili{Lio}
and Palu{\Q}e \it{ləma} ‘tongue’ (and similar forms in many languages
throughout the \tsc{\ili{Bima}-Lembata} subgroup) are more easily linked
to PMP *həma ‘tongue than to PCEMP **maya ‘tongue’.\footnote{
		Similar forms for ‘tongue’ with an \it{l}-onset seems to be fairly widely
		distributed, including in over 40 languages in Vanuatu.
		However, Malcolm Ross (pers.
		comm.) believes that the forms in Vanuatu have a separate explanation
		in their historical phonologies.
		We are simply observing their similarities of form
		and meaning.}
Indeed, it is not always straightforward to determine if a particular
form for ‘tongue’ is a derivative of **maya or *həma.
Thus, for instance, we have identified \ili{Taliabu} \it{maa-n} ‘tongue’
as reflex of **maya (with irregular loss of medial **y [see \srf{08-01-02}]),
but we note that the restructuring of PMP *əpat > \it{ŋ-haaʔ}
‘four’ and *ənəm > \it{nooŋ} ‘six’
would also parallel *həma > \it{maa-n}.

\begin{table}[h]
	\centering\tcp{Etyma for `tongue'}{15-08a}
	\st{6pt}\begin{tabular}{lll}\lsptoprule
	Level &	Etymon	&	note \\	\midrule
	\tsc{PAn}	&	*Səma	&	(PMP *həma)\\
	\tsc{PAn}	&	*lidam	&	\\
	PMP	&	*dilaq	&	\\
	PMP?	&	*ləlaw	&	(see discussion above)\\
	P(W)MP	&	*zəlaq	&	\\
	PCEMP	&	**maya	&	(perhaps ultimately derived\\
				&					&	from *həma? how? evidence?)\\
	\tsc{POc}	&	**maya	&	\\
	Proto-Polynesian	&	*qalelo	&	(possibly related to *ləlaw above)\\
		\lspbottomrule
	\end{tabular}
\end{table}

Possible (but irregular) reflexes of PMP *dilaq ‘tongue,
lick’ may also be found in the region with lowering of *i (or
more likely related to P(W)MP *zəlaq ‘tongue’,\footnote{
		Forms such as \it{elo} and \it{lelo} are extremely common
		in languages of south\ili{east} Sulawesi (cf.
		\citealp[105]{Mead1999}), and are assigned to PMP *dilaq ‘tongue,
		lick’ in \citet{GreenhillGray2008}.}
or an alternate competing form
such as *ləlaw ‘tongue’),
such as \ili{Bima} \it{lera} ‘tongue’ (\tsc{\ili{Bima}-Lembata}); \ili{Liana-Seti}
\it{lela-} ‘tongue’ (\tsc{\ili{Ambon}-Seram}), \ili{Masiwang} \it{lehaʔ} (*l > \it{l,
h}; STB), \ili{Kowiai} \it{yera} ‘tongue’ (\ili{Kowiai} *l > \it{r} is regular;
STB); SHWNG \ili{Taba} \it{plelo},
\ili{Buli} \it{papleo}, \ili{Sawai} \it{pɛplɛ-o},
\ili{Biak} \it{kaprer}, all the way out to Proto-Polynesian *qalelo ‘tongue’.
So consequently the story of ‘tongue’ \ili{east} of the Blust line
is far more complex than simply **maya ‘tongue’ being a replacement
innovation for PMP *dilaq in the CEMP region \citep[740]{Blust2013}.

In addition to the proto forms in example \trf{15-08a} there are
additional reconstructions not listed here.
The words shown in \trf{15-08} support **ləma.

\begin{table}[p]
	\centering\tcp{Data reflecting **ləma ‘tongue’}{15-08}
	\st{6pt}\begin{tabular}{l|ll}\lsptoprule
Language	&	Form	&	Node / Comment\\	\midrule
	&	**ləma	&	possibly related to PMP *həma\\
\ili{Soboyo}	&	{\hr\hr}dama-in	&	\tsc{Cel}, \tsc{Tal}\\
\ili{Kadai}	&	{\hr\hr}lamaa	&	\tsc{Cel}, \tsc{Tal}\\	\midrule
\ili{Manggarai}	&	{\hr\hr}ləma	&	\tsc{BL}, \tsc{W. Flores}\\
\ili{Ngadha}	&	{\hr\hr}ləma	&	\tsc{BL}, \tsc{C. Flores}\\
So{\Q}a	&	{\hr\hr}ləma	&	\tsc{BL}, \tsc{C. Flores}\\
\ili{Nage}	&	{\hr\hr}ləma	&	\tsc{BL}, \tsc{C. Flores}\\
\ili{Ende}	&	{\hr\hr}ɹəma	&	\tsc{BL}, \tsc{C. Flores}\\
Li'o	&	{\hr\hr}ləma	&	\tsc{BL}, \tsc{C. Flores}\\
Palu{\Q}e	&	{\hr\hr}ləma	&	\tsc{BL}, \tsc{C. Flores}\\
\ili{Kambera}	&	{\hr\hr}lɐma	&	\tsc{BL}, \tsc{Sumba}\\
\ili{Kodi}	&	{\hr\hr}lamːa	&	\tsc{BL}, \tsc{Sumba}\\
\ili{Wanukaka}	&	{\hr\hr}lamːa	&	\tsc{BL}, \tsc{Sumba}\\
\ili{Anakalang}	&	{\hr\hr}lɐma	&	\tsc{BL}, \tsc{Sumba}\\
\ili{Laboya}	&	{\hr\hr}lama	&	\tsc{BL}, \tsc{Sumba}\\
\ili{Mamboru}	&	{\hr\hr}lɐma	&	\tsc{BL}, \tsc{Sumba}\\
\ili{Weyewa}	&	{\hr\hr}lomːa	&	\tsc{BL}, \tsc{Sumba}\\	\midrule
\ili{Kemak}	&	{\hr\hr}lama-r	&	\tsc{CT}\\
S. \ili{Mambae}	&	{\hr\hr}lama {\tl} laam	&	\tsc{CT}\\
\ili{Tokodede}	&	{\hr\hr}lama	&	\tsc{CT}\\	\midrule
West \ili{Damer}	&	{\hr\hr}lelɛmma	&	\tsc{TB}, \tsc{SW. Maluku}\\
Proto-\ili{Babar}	&	**nama	&	\tsc{TB}, \tsc{SW. Maluku}; irregular *l > **n\\
\ili{Kisar}	&	{\hr\hr}nama	&	\tsc{TB}, \tsc{SW. Maluku}; irregular *l > \it{n}\\
\ili{Luang}	&	{\hr\hr}namə	&	\tsc{TB}, \tsc{SW. Maluku}; irregular *l > \it{n}\\	\midrule
\ili{Batumerah}	&	{\hr\hr}lama-a	&	\tsc{AS}\\
\ili{Luhu}	&	{\hr\hr}lamaa	&	\tsc{AS}\\	\midrule
\ili{Watubela}	&	{\hr\hr}lama	&	\tsc{STB}\\	\midrule
Vanuatu lgs.	&	{\hr\hr}ləme-, leme-, lama-	&	40+ languages; possibly unrelated\\
		\lspbottomrule
	\end{tabular}
\end{table}

Do the data in \trf{15-08} have value for subgrouping at a higher
level, or only at lower levels? How does one explain the strong but
disjointed pockets in \tsc{\ili{Bima}-Lembata},
\tsc{\ili{Central} Timor}, and some of the \tsc{SW.
Maluku} node of \tsc{Timor-Babar}? How does one explain \tsc{Celebic}
languages \ili{Soboyo}-\ili{Kadai}, and others such as \ili{Luhu}
and \ili{Watubela} standing off by themselves in Maluku?